\subsectionaddtolist{الف}
\[
    L = \{w \in \{a\}^* : |w| \  mod \ 3 > 0\}
    \]
این زبان شامل تمام رشته‌هایی از حرف $a$ است که طول آن‌ها مضرب ۳ نباشد. به عبارت دیگر، باقی‌مانده طول رشته بر ۳ برابر با ۱ یا ۲ باشد. 

سه متغیر $S$ (باقی‌مانده ۰)، $A$ (باقی‌مانده ۱) و $B$ (باقی‌مانده ۲) را تعریف می‌کنیم. از آنجا که باقی‌مانده‌های ۱ و ۲ مورد قبول هستند، متغیرهای $A$ و $B$ می‌توانند به رشته تهی ($\lambda$) ختم شوند.

\[
\begin{aligned}
    S &\rightarrow aA \\
    A &\rightarrow aB \mid \lambda \\
    B &\rightarrow aS \mid \lambda
\end{aligned}
\]

\subsectionaddtolist{ب}
\[
    L = \{a^n b^m : n + m = odd\}
    \]
رشته‌ها از تعدادی $a$ و سپس تعدادی $b$ تشکیل شده‌اند، به طوری که مجموع طول کل رشته فرد باشد. این یعنی زوجیت تعداد $a$ها و $b$ها متفاوت است.

\begin{itemize}
    \item $S$: حالت شروع (تعداد $a$ زوج، تولید $b$ شروع نشده)
    \item $A$: تعداد $a$ فرد، تولید $b$ شروع نشده
    \item $B_e$: مرحله تولید $b$، به طوری که مجموع کل طول تا این لحظه زوج است
    \item $B_o$: مرحله تولید $b$، به طوری که مجموع کل طول تا این لحظه فرد است
\end{itemize}

\[
\begin{aligned}
    S &\rightarrow aA \mid bB_o \\
    A &\rightarrow aS \mid bB_e \mid \lambda \\
    B_e &\rightarrow bB_o \\
    B_o &\rightarrow bB_e \mid \lambda
\end{aligned}
\]

\subsectionaddtolist{ج}

\[
L = \{vwv : v, w \in \{a,b\}^*, |v| = 2\}
\]

رشته با یک زیررشته به طول ۲ ($v$) شروع شده و با همان زیررشته به پایان می‌رسد. از آنجا که $v \in \{aa, ab, ba, bb\}$ است، گرامر مسیرهای تولید را بر اساس ابتدای رشته تفکیک کرده و پایانی متناظر با آن را تضمین می‌کند.


\[
\begin{aligned}
    S &\rightarrow aS_{a} \mid bS_{b} \\
    S_{a} &\rightarrow aS_{aa} \mid bS_{ab} \\
    S_{b} &\rightarrow aS_{ba} \mid bS_{bb} \\
    \\
    S_{aa} &\rightarrow aS_{aa} \mid bS_{aa} \mid aA_{aa} \\
    A_{aa} &\rightarrow a \\
    \\
    S_{ab} &\rightarrow aS_{ab} \mid bS_{ab} \mid aA_{ab} \\
    A_{ab} &\rightarrow b \\
    S_{ba} &\rightarrow aS_{ba} \mid bS_{ba} \mid bA_{ba} \\
    A_{ba} &\rightarrow a \\
    \\
    S_{bb} &\rightarrow aS_{bb} \mid bS_{bb} \mid bA_{bb} \\
    A_{bb} &\rightarrow b
\end{aligned}
\]
% ----------------------------------------------------------------------
\subsectionaddtolist{د}

\[
\text{ مجموعه تمام اعداد در مبنای ۴ که به عدد ۷ بخش‌پذیرند}
\]

الفبا $\Sigma = \{0, 1, 2, 3\}$ است. متغیرهای $R_0$ تا $R_6$ نشان‌دهنده باقی‌مانده عدد تولید شده تا این لحظه بر ۷ هستند. با اضافه شدن رقم $d$ در مبنای ۴، باقی‌مانده جدید از رابطه $R_{\text{new}} = (R_{\text{old}} \times 4 + d) \pmod 7$ محاسبه می‌شود.

با محاسبه انتقال‌ها برای تمامی حالات، گرامر زیر حاصل می‌شود:

\[
\begin{aligned}
    R_0 &\rightarrow 0R_0 \mid 1R_1 \mid 2R_2 \mid 3R_3 \mid \lambda \\
    R_1 &\rightarrow 0R_4 \mid 1R_5 \mid 2R_6 \mid 3R_0 \\
    R_2 &\rightarrow 0R_1 \mid 1R_2 \mid 2R_3 \mid 3R_4 \\
    R_3 &\rightarrow 0R_5 \mid 1R_6 \mid 2R_0 \mid 3R_1 \\
    R_4 &\rightarrow 0R_2 \mid 1R_3 \mid 2R_4 \mid 3R_5 \\
    R_5 &\rightarrow 0R_6 \mid 1R_0 \mid 2R_1 \mid 3R_2 \\
    R_6 &\rightarrow 0R_3 \mid 1R_4 \mid 2R_5 \mid 3R_6
\end{aligned}
\]
